\documentclass[10pt,twocolumn]{article}

\usepackage[T1]{fontenc}
\usepackage{lmodern}
\usepackage[margin=0.56in]{geometry}
\usepackage{amsmath,amssymb,mathtools,bm}
\usepackage{xcolor}
\usepackage{enumitem}
\usepackage{booktabs}
\usepackage[hidelinks]{hyperref}

\setlength{\columnsep}{0.24in}
\setlength{\parindent}{0pt}
\setlength{\parskip}{3pt}
\setlist[itemize]{leftmargin=1.15em,itemsep=1pt,topsep=2pt}
\setlist[enumerate]{leftmargin=1.25em,itemsep=1pt,topsep=2pt}
\allowdisplaybreaks

\newcommand{\ii}{\mathrm{i}}
\newcommand{\ket}[1]{\left|#1\right\rangle}
\newcommand{\bra}[1]{\left\langle#1\right|}
\newcommand{\braket}[2]{\left\langle #1\,\middle|\,#2\right\rangle}
\newcommand{\norm}[1]{\left\lVert #1\right\rVert}
\newcommand{\ip}[2]{\left\langle #1,#2\right\rangle}
\newcommand{\R}{\mathbb R}
\newcommand{\C}{\mathbb C}
\newcommand{\Hh}{\mathcal H}
\newcommand{\argmin}{\operatorname*{arg\,min}}
\newcommand{\argmax}{\operatorname*{arg\,max}}
\newcommand{\diag}{\operatorname{diag}}
\newcommand{\rank}{\operatorname{rank}}
\newcommand{\spanop}{\operatorname{span}}
\newcommand{\clip}{\operatorname{clip}}
\newcommand{\Real}{\operatorname{Re}}
\newcommand{\Imag}{\operatorname{Im}}
\newcommand{\Id}{I}
\newcommand{\eps}{\varepsilon}
\newcommand{\dd}{\mathrm d}
\newcommand{\transpose}{\mathsf T}
\newcommand{\adj}{\dagger}

\definecolor{boxgray}{RGB}{246,246,243}
\definecolor{linegray}{RGB}{213,213,205}
\definecolor{deepblue}{RGB}{20,70,125}
\definecolor{deepred}{RGB}{130,35,35}
\definecolor{deepgreen}{RGB}{35,95,60}
\definecolor{deepviolet}{RGB}{85,55,125}

\newcommand{\micro}[1]{%
  \par\smallskip
  \noindent\textcolor{deepblue}{\textbf{Micro-intuition:}} #1
  \par\smallskip}

\newcommand{\deriv}{%
  \par\noindent\textcolor{deepviolet}{\textbf{Mathematical derivation:}}\par\vspace{-1mm}}

\newcommand{\phys}[1]{%
  \par\smallskip
  \noindent\textcolor{deepgreen}{\textbf{Physical interpretation:}} #1
  \par\smallskip}

\newcommand{\geom}[1]{%
  \par\smallskip
  \noindent\textcolor{deepviolet}{\textbf{Geometric picture:}} #1
  \par\smallskip}

\newcommand{\caution}[1]{%
  \par\smallskip
  \noindent\fcolorbox{linegray}{boxgray}{%
  \begin{minipage}{0.965\columnwidth}
  \textcolor{deepred}{\textbf{Caution:}} #1
  \end{minipage}}
  \par\smallskip}

\newcommand{\paperbox}[1]{%
  \par\smallskip
  \noindent\fcolorbox{linegray}{boxgray}{%
  \begin{minipage}{0.965\columnwidth}
  \textbf{Paper-compatible translation:} #1
  \end{minipage}}
  \par\smallskip}

\newcommand{\implbox}[1]{%
  \par\smallskip
  \noindent\fcolorbox{linegray}{boxgray}{%
  \begin{minipage}{0.965\columnwidth}
  \textbf{Implementation-level reading:} #1
  \end{minipage}}
  \par\smallskip}

\title{\vspace{-1.1cm}AP--McLachlan Time Dynamics\\
\large A Pedagogical Companion for the Inverse, Repair, Append, and Prune Logic}
\author{}
\date{\vspace{-0.95cm}}

\begin{document}
\maketitle
\vspace{-0.9cm}

\section*{0. Local convention}

\micro{The whole note should be read at one checkpoint unless a previous or
future checkpoint is written explicitly.}

\deriv
Let
\[
\begin{aligned}
\Hh&\cong \C^d,\\
\ket{\psi}&\in\Hh,\qquad \braket{\psi}{\psi}=1,\\
\theta&=(\theta_j)_{j=1}^{n}\in\R^n,\\
U(\theta)&:\Hh\to\Hh,\qquad U(\theta)^\dagger U(\theta)=\Id,\\
\ket{\psi(\theta)}&=U(\theta)\ket{\psi_{\rm ref}}\in\Hh.
\end{aligned}
\]
At the checkpoint, write \(\ket{\psi}=\ket{\psi(\theta)}\).

\phys{The ansatz state is a vector in the Hilbert space. The parameters are
classical real coordinates. The variational dynamics tries to move \(\theta\)
so that the state velocity follows the Schrödinger velocity.}

\geom{Picture \(\R^n\) as the coordinate chart and the variational states as a
curved surface inside Hilbert space. A parameter-space arrow \(v\) becomes a
state-space arrow \(Tv\) only after the tangent map pushes it onto that surface.
McLachlan asks which pushed arrow lies closest to the target drift \(\ket b\).}

\paperbox{Paper II often writes the checkpoint as \(k\), for example
\(\ket{\psi_k}=U(\theta_k)\ket{\psi_{\rm ref}}\). This companion suppresses \(k\) until a
time comparison is actually needed.}

\section*{1. The target velocity}

\micro{The exact Schrödinger velocity is projected horizontally so global phase
does not become a fake variational direction.}

\deriv
Let \(H(t):\Hh\to\Hh\) be Hermitian and let
\[
\begin{aligned}
e_\psi&=\bra{\psi}H(t)\ket{\psi}\in\R,\\
\Pi_\psi&=\Id-\ket{\psi}\bra{\psi}:\Hh\to\Hh,\\
\ket{b}&=-\ii\bigl(H(t)-e_\psi\Id\bigr)\ket{\psi}\in\Hh.
\end{aligned}
\]
Then
\[
\begin{aligned}
\braket{\psi}{b}
&=-\ii\bra{\psi}H(t)\ket{\psi}
+\ii e_\psi\braket{\psi}{\psi}\\
&=-\ii e_\psi+\ii e_\psi\\
&=0,
\end{aligned}
\]
and
\[
\Pi_\psi\ket{b}=\ket{b}.
\]

\phys{The ket \(\ket b\) is the physically relevant instantaneous drift. The
removed component is parallel to \(\ket\psi\), so it changes only global phase.}

\micro{The norm of the projected drift is the local size of the exact motion
that the tangent space is asked to represent.}

\deriv
\[
\begin{aligned}
\norm{\ket b}^2
&=\braket{b}{b}\\
&=\bra{\psi}\bigl(H(t)-e_\psi\Id\bigr)^2\ket{\psi}\\
&=\bra{\psi}H(t)^2\ket{\psi}-e_\psi^2\\
&=\operatorname{Var}_{\psi}(H(t)).
\end{aligned}
\]

\phys{A small \(\norm{\ket b}^2\) means the exact horizontal velocity is small at that
checkpoint. A large residual fraction is physically less alarming when the
absolute drift itself is tiny, although the current noiseless diagnostics can
still track the fraction.}

\section*{2. Tangents from the ansatz}

\micro{The tangent matrix is the linear map from parameter velocity to state
velocity.}

\deriv
For \(1\le j\le n\),
\[
\begin{aligned}
\partial_j\ket{\psi}
&=\frac{\partial}{\partial\theta_j}\ket{\psi(\theta)}\in\Hh,\\
\ket{t_j}&=\Pi_\psi\,\partial_j\ket{\psi}\in\Hh.
\end{aligned}
\]
Define the tangent matrix
\[
T=
\begin{bmatrix}
| & & |\\
\ket{t_1}&\cdots&\ket{t_n}\\
| & & |
\end{bmatrix}
\in\C^{d\times n}.
\]
For \(v=(v_j)_{j=1}^n\in\R^n\),
\[
\begin{aligned}
Tv
&=\sum_{j=1}^n v_j\ket{t_j}\in\Hh,\\
(Tv)_a
&=\sum_{j=1}^n T_{aj}v_j,\qquad 1\le a\le d.
\end{aligned}
\]

\phys{A real vector \(v\) is a candidate parameter velocity. Multiplying by
\(T\) converts parameter velocity into the represented horizontal state
velocity.}

\geom{The columns of \(T\) are the local arrows drawn by the ansatz coordinates.
If two columns draw nearly the same arrow, the coordinate chart has redundancy.
If one column is very short, that coordinate barely moves the state.}

\micro{For a product ansatz, noncommuting operator order matters in the tangent
columns.}

\deriv
Let
\[
\begin{aligned}
U(\theta)
&=U_n(\theta_n)\cdots U_1(\theta_1),\\
U_j(\theta_j)
&=\exp(-\ii\theta_j A_j),\qquad A_j=A_j^\dagger.
\end{aligned}
\]
Then
\[
\begin{aligned}
\partial_j U(\theta)
&=
U_n\cdots U_{j+1}(\partial_jU_j)U_{j-1}\cdots U_1,\\
\partial_jU_j
&=-\ii A_jU_j,\\
\partial_j\ket{\psi}
&=
U_n\cdots U_{j+1}(-\ii A_j)U_j\cdots U_1\ket{\psi_{\rm ref}}.
\end{aligned}
\]
Thus
\[
\ket{t_j}
=\Pi_\psi U_n\cdots U_{j+1}(-\ii A_j)U_j\cdots U_1\ket{\psi_{\rm ref}}.
\]

\phys{The \(j\)-th tangent is not simply \(-\ii A_j\ket\psi\) unless the generator is
the last applied factor or the implementation has already pushed the generator
into the correct dressed frame. The code's executor matters because it fixes
this ordered product.}

\implbox{The executor is the concrete rule that maps the ordered ansatz record
and \(\theta\) to \(U(\theta)\ket{\psi_{\rm ref}}\), and to the tangent columns that
feed \(T\). Loading a final state can bypass state preparation, but McLachlan
still needs the tangent map of the ansatz if parameters are propagated.}

\section*{3. The McLachlan objective}

\micro{The least-squares problem chooses the parameter velocity whose tangent
velocity best approximates \(\ket b\).}

\deriv
Let \(\dot\theta\in\R^n\) be a candidate parameter velocity. Using the complex
Hilbert-space inner product,
\[
\begin{alignedat}{2}
m(\dot\theta)
&=\norm{T\dot\theta-\ket b}^2
&&\quad\scriptstyle(\text{definition})\\
&=(T\dot\theta-\ket b)^\dagger(T\dot\theta-\ket b)
&&\quad\scriptstyle(\norm{x}^2=x^\dagger x)\\
&=\dot\theta^\top T^\dagger T\dot\theta
-\dot\theta^\top T^\dagger\ket b
-\bra b T\dot\theta
+\braket{b}{b}
&&\quad\scriptstyle(\text{expand})\\
&=\dot\theta^\top T^\dagger T\dot\theta
-2\Real\!\left(\dot\theta^\top T^\dagger\ket b\right)
&&\quad\scriptscriptstyle(z+z^*=2\Real z)\\
&\quad+\norm{\ket b}^2\\
&=\dot\theta^\top\Real(T^\dagger T)\dot\theta
-2\dot\theta^\top\Real(T^\dagger\ket b)
&&\quad\scriptscriptstyle(\dot\theta\in\R^n)\\
&\quad+\norm{\ket b}^2.
\end{alignedat}
\]
Here the fourth line uses
\(\bra bT\dot\theta=(\dot\theta^\top T^\dagger\ket b)^*\). The last line uses
the fact that the imaginary part of the Hermitian quadratic form does not
contribute for real \(\dot\theta\). With
\[
G=\Real(T^\dagger T),
\qquad
f=\Real(T^\dagger\ket b),
\]
the objective is
\[
m(\dot\theta)
=\dot\theta^\top G\dot\theta-2\dot\theta^\top f+\norm{\ket b}^2.
\]

\micro{The implementation form is the same normal equation written entrywise.}

\deriv
Let
\[
\begin{aligned}
T&=(T_{aj})_{1\le a\le d,\;1\le j\le n}\in\C^{d\times n},\\
\ket b&\leftrightarrow b=(b_a)_{a=1}^{d}\in\C^d,\\
\dot\theta&=(\dot\theta_j)_{j=1}^{n}\in\R^n.
\end{aligned}
\]
The definitions
\[
G=\Real(T^\dagger T),
\qquad
f=\Real(T^\dagger\ket b)
\]
mean
\[
\begin{aligned}
G_{ij}
&=\Real\sum_{a=1}^{d}T_{ai}^*T_{aj},
\qquad
1\le i,j\le n,\\
f_i
&=\Real\sum_{a=1}^{d}T_{ai}^*b_a,
\qquad
1\le i\le n.
\end{aligned}
\]
Thus \(G\dot\theta=f\) is the component equation
\[
\sum_{j=1}^{n}G_{ij}\dot\theta_j=f_i,
\qquad
1\le i\le n,
\]
or, after substituting the primitive overlap definitions,
\[
\sum_{j=1}^{n}
\left(
\Real\sum_{a=1}^{d}T_{ai}^*T_{aj}
\right)
\dot\theta_j
=
\Real\sum_{a=1}^{d}T_{ai}^*b_a,
\qquad
1\le i\le n.
\]
The component derivative uses dummy-index relabeling and symmetry:
\[
\begin{aligned}
\frac{\partial}{\partial\dot\theta_i}
\left(\sum_{p,q}\dot\theta_pG_{pq}\dot\theta_q\right)
&=
\sum_q G_{iq}\dot\theta_q+\sum_p\dot\theta_pG_{pi}\\
&=
\sum_q G_{iq}\dot\theta_q+\sum_q\dot\theta_qG_{qi}\\
&=
2\sum_qG_{iq}\dot\theta_q,
\end{aligned}
\]
where \(G_{qi}=G_{iq}\). Thus the \(i\)-th stationarity equation is
\[
2\sum_qG_{iq}\dot\theta_q-2f_i=0.
\]

\phys{The real part enters because the variational parameters are real, so the
complex Hilbert-space overlaps become a real quadratic problem in the parameter
velocity \(\dot\theta\). Each \(G_{ij}\) is a tangent--tangent overlap, each
\(f_i\) is a tangent--drift overlap, and the implementation fills these real
arrays before solving for the coordinate velocity.}

\section*{4. Why \(G\dot\theta=f\)}

\micro{A normal equation is the gradient-zero system for least squares:
\(\min_x\norm{Ax-y}^2\) gives \(A^\dagger Ax=A^\dagger y\). Here \(A=T\),
\(x=\dot\theta\), \(y=\ket b\), and real coordinates give \(G\dot\theta=f\).}

\deriv
\[
\begin{aligned}
m(\dot\theta)&=\dot\theta^\top G\dot\theta-2\dot\theta^\top f+\norm{\ket b}^2,\\
G^\top&=G,\qquad G\succeq0.
\end{aligned}
\]
Using the vector-gradient identities
\[
\begin{aligned}
\nabla_{\dot\theta}(\dot\theta^\top G\dot\theta)&=2G\dot\theta,\\
\nabla_{\dot\theta}(2\dot\theta^\top f)&=2f,
\end{aligned}
\]
Thus
\[
\nabla_{\dot\theta}m(\dot\theta)=2G\dot\theta-2f.
\]
The stationary velocity satisfies
\[
G\dot\theta=f.
\]

\phys{The McLachlan velocity is the parameter velocity where an infinitesimal
coordinate change cannot further reduce the state-space mismatch.}

\micro{Non-identifiability means distinct coordinate velocities create the same
first-order state velocity.}

\deriv
Two coordinate velocities are locally indistinguishable when their difference
lies in the nullspace of the tangent map:
\[
T(\dot\theta+a)=T\dot\theta
\quad\Longleftrightarrow\quad
Ta=0.
\]
Equivalently,
\[
\dot\theta\sim\dot\theta+a
\quad\text{when}\quad
a\in\operatorname{null}T.
\]
The identifiable object is the represented velocity, not the raw coordinate
velocity:
\[
[\,\dot\theta\,]\longmapsto T\dot\theta,
\qquad
[\,\dot\theta\,]=\dot\theta+\operatorname{null}T.
\]

\phys{A non-identifiable coordinate direction is locally gauge-like: moving in
that coordinate changes the parameter vector but not the quantum state to first
order. McLachlan must solve on the visible quotient coordinates, not on every
raw coordinate direction as though all directions were physical.}

\micro{Singular \(G\) means some nonzero coordinate direction is invisible to
the tangent map.}

\deriv
A square matrix is singular when it has a nonzero null direction. A null
direction is a coordinate vector that the matrix sends to zero:
\[
\exists\,a\in\R^n,\qquad a\ne0,\qquad Ga=0.
\]
For the McLachlan Gram matrix,
\[
\begin{aligned}
G&=\Real(T^\dagger T),\\
a^\top Ga
&=a^\top\Real(T^\dagger T)a\\
&=\Real\bigl((Ta)^\dagger(Ta)\bigr)\\
&=\norm{Ta}^2.
\end{aligned}
\]
Thus, in the exact positive-semidefinite geometry,
\[
Ga=0
\quad\Longrightarrow\quad
a^\top Ga=0
\quad\Longrightarrow\quad
Ta=0.
\]

\phys{The coordinate direction \(a\) changes parameters, but it produces no
first-order state velocity. The ansatz has a redundant or ineffective local
parameter direction.}

\geom{A null direction is a coordinate arrow that vanishes when projected onto
the variational surface. Moving along that coordinate can look large in
parameter space while producing no visible first-order motion of the state.}

\micro{The pseudoinverse solves only on the range of \(G\), and sets null
directions to zero.}

\deriv
The nullspace is the set of all null directions. The range is the set of all
vectors the matrix can output:
\[
\operatorname{null}G=\{a\in\R^n:Ga=0\},
\qquad
\operatorname{ran}G=\{Gy:y\in\R^n\}.
\]
The theorem used here is the fundamental subspace theorem for a real symmetric
matrix: the range is the orthogonal complement of the nullspace,
\[
\operatorname{ran}G=(\operatorname{null}G)^\perp.
\]
For symmetric positive-semidefinite \(G\), this gives the orthogonal direct
sum
\[
\R^n=\operatorname{ran}G\oplus\operatorname{null}G.
\]
The symbol \(\oplus\) means each vector splits uniquely into a range component
plus a nullspace component.
The Moore--Penrose pseudoinverse theorem says that \(GG^+\) is the orthogonal
projector onto \(\operatorname{ran}G\), and that \(G^+f\) is the
minimum-norm least-squares solution of \(Gx\approx f\).
If \(f\in\operatorname{ran}G\), then \(G\dot\theta=f\) has supported solutions.
They differ by null directions:
\[
\dot\theta=G^+f+n,
\qquad
n\in\operatorname{null}G.
\]
The minimum-norm supported solution chooses \(n=0\):
\[
\dot\theta=G^+f.
\]
Then
\[
G\dot\theta=GG^+f=f.
\]
If \(f\notin\operatorname{ran}G\), the full equation \(G\dot\theta=f\) has no
exact solution. The pseudoinverse returns the least-squares supported solve:
\[
\dot\theta=G^+f,
\qquad
G\dot\theta=GG^+f=P_{\operatorname{ran}G}f.
\]
The unsupported force component is
\[
f_\perp=f-P_{\operatorname{ran}G}f.
\]

\phys{The retained tangent coordinates respond only to the component of the
force \(f\) that lies in the supported metric range. The nullspace component
would create parameter motion without state-space motion, so the minimum-norm
solve removes it.}

\geom{The range is the visible subspace of tangent forces. The projector
\(P_{\operatorname{ran}G}\) drops the part of \(f\) that points outside that
visible subspace, leaving only the force shadow the current tangent chart can
answer.}

\micro{A near-null direction makes the ordinary inverse large because the solve
divides by the metric strength of that direction.}

\deriv
The theorem used here is the real spectral theorem: every real symmetric matrix
has an orthonormal eigenbasis. Thus, for symmetric positive-semidefinite \(G\),
there are orthonormal vectors \(v_a\) and nonnegative eigenvalues \(s_a\) such
that
\[
Gv_a=s_av_a,\qquad v_a^\top v_a=1,\qquad s_a\ge0.
\]
Because the vectors \(v_a\) form an orthonormal basis, every vector in
\(\R^n\), including \(f\) and \(\dot\theta\), has a unique expansion in these
eigenmodes.
The number \(s_a\) is the physical tangent strength of the parameter direction
\(v_a\):
\[
s_a=v_a^\top Gv_a=\norm{Tv_a}^2.
\]
Expand the force and velocity in the same eigenbasis:
\[
f=\sum_a f_av_a,\qquad
\dot\theta=\sum_a y_av_a,\qquad
f_a=v_a^\top f.
\]
The normal equation \(G\dot\theta=f\) becomes, mode by mode,
\[
s_ay_a=f_a.
\]
Therefore an ordinary inverse would require
\[
y_a=\frac{f_a}{s_a}.
\]
This is the precise instability:
\[
s_a\to0^+,\quad f_a\ne0
\quad\Longrightarrow\quad
|y_a|=\frac{|f_a|}{s_a}\to\infty.
\]
When \(G\) is nonsingular, the inverse response map is
\[
\Phi_G:\R^n\to\R^n,
\qquad
\Phi_G(\delta f)=G^{-1}\delta f.
\]
For a perturbation \(\delta f=\sum_a \delta f_av_a\),
\[
\Phi_G(\delta f)
=
\sum_a\frac{\delta f_a}{s_a}v_a.
\]
The operator norm of this response map is
\[
\norm{\Phi_G}_{2\to2}
=
\sup_{\delta f\ne0}
\frac{\norm{G^{-1}\delta f}}{\norm{\delta f}}
=
\frac{1}{s_{\min}},
\qquad
s_{\min}=\min_a s_a.
\]
The supremum is attained when \(\delta f\) points along an eigenmode with
eigenvalue \(s_{\min}\). The spectral condition number is
\[
\kappa(G)=\frac{s_{\max}}{s_{\min}},
\qquad
s_{\max}=\max_a s_a,
\]
so a large \(\kappa(G)\) means the weakest identifiable mode is much weaker
than the strongest one. A small force error along the weakest eigenmode can
create a large coordinate-velocity error.
If \(s_a=0\) exactly, then \(1/s_a\) is undefined. A supported inverse avoids
this by using a retention indicator \(\chi_a\):
\[
y_a^{\rm supp}
=\chi_a\frac{f_a}{s_a}.
\]
Here \(\chi_a=0\) on discarded null or near-null modes.
A damped supported inverse further replaces \(s_a\) by
\(s_a+\epsilon\) on retained modes.

\phys{A near-null coordinate direction barely moves the quantum state. If a
finite force component or a small force error is assigned to that weak
direction, an ordinary inverse compensates by making the parameter velocity
large. Truncation and damping prevent the algorithm from turning weakly
identifiable coordinate motion into an unstable time step.}

\geom{In the eigenbasis, \(G\) turns the coordinate chart into principal axes.
The axis \(v_a\) has visible tangent length \(s_a^{1/2}\). A near-flat axis
requires a long coordinate arrow \(y_av_a\) to make a small state-space arrow,
which is why \(y_a=f_a/s_a\) becomes unstable when \(s_a\) is tiny.}

\section*{5. Explained drift}

\micro{The quadratic form \(f^\top G^+f\) is the amount of squared drift
explained after the best tangent coefficients have been optimized out.}

\deriv
Let
\[
\dot\theta=G^+f.
\]
The represented velocity is
\[
\ket u=T\dot\theta\in\Hh.
\]
Its squared norm is
\[
\begin{aligned}
\norm{\ket u}^2
&=(T\dot\theta)^\dagger(T\dot\theta)\\
&=\dot\theta^\top T^\dagger T\dot\theta.
\end{aligned}
\]
Since \(\dot\theta\in\R^n\),
\[
\dot\theta^\top T^\dagger T\dot\theta
=\dot\theta^\top\Real(T^\dagger T)\dot\theta
=\dot\theta^\top G\dot\theta.
\]
Substitute \(\dot\theta=G^+f\):
\[
\norm{\ket u}^2
=f^\top G^+GG^+f.
\]
For a symmetric positive semidefinite \(G\),
\[
G^+GG^+=G^+.
\]
Thus
\[
\boxed{\norm{T\dot\theta}^2=f^\top G^+f.}
\]

\phys{The overlap vector \(f\) by itself is not an explained length. The metric
correction \(G^+\) converts raw shadows into nonorthogonal coefficients, then
the quadratic form reports the length of the fitted tangent velocity.}

\micro{The same expression appears by substituting the optimizer back into the
least-squares residual.}

\deriv
\[
\begin{aligned}
m(\dot\theta)
&=\dot\theta^\top G\dot\theta-2\dot\theta^\top f+\norm{\ket b}^2\\
&=f^\top G^+GG^+f-2f^\top G^+f+\norm{\ket b}^2\\
&=\norm{\ket b}^2-f^\top G^+f.
\end{aligned}
\]
Therefore
\[
\boxed{
\min_{v\in\R^n}\norm{Tv-\ket b}^2
=
\norm{\ket b}^2-f^\top G^+f.
}
\]

\phys{The velocity variable disappeared because it was optimized out. The
reported residual is the minimized error left after the tangent span made its
best local attempt to reconstruct \(\ket b\).}

\section*{6. Structural tangent miss}

\micro{The structural miss asks whether the current tangent span can explain
the Schrödinger drift after the best supported fit.}

\deriv
Use the same supported pseudoinverse convention as the projection geometry;
Section 7 expands this inverse spectrally. The current structural miss is
\[
\rho(T,\ket b)
=
\frac{
\left[
\norm{\ket b}^2-f^\top G^+f
\right]_+
}{\norm{\ket b}^2}
\]
and equivalently
\[
\rho(T,\ket b)
=
\frac{
\left[
\min_{v\in\R^n}\norm{Tv-\ket b}^2
\right]_+
}{\norm{\ket b}^2}.
\]
The first expression is the computable Gram form; the second expression is the
least-squares meaning.

\phys{Large \(\rho(T,\ket b)\) means the manifold is missing physical tangent
directions. This is the structural reason to consider append.}

\caution{The denominator can have a small numerical floor in the implementation.
The floor prevents division by a tiny drift norm; it is not a new physical
quantity.}

\section*{7. Supported inverse}

\micro{The supported inverse keeps retained eigenmodes and truncates untrusted
eigenmodes.}

\deriv
By the real spectral theorem, write the symmetric positive-semidefinite Gram
matrix as
\[
G=V\diag(s_1,\ldots,s_n)V^\top,
\qquad
s_a\ge0,
\qquad
V^\top V=\Id.
\]
The columns of \(V\) are the orthonormal eigenmodes \(v_a\), and the numbers
\(s_a\) are their metric strengths.
Let
\[
s_{\max}=\max_a s_a,
\qquad
\chi_a=\mathbf 1[s_a\ge\tau s_{\max}].
\]
A retained mode has \(\chi_a=1\). A truncated mode has \(\chi_a=0\), meaning its
inverse weight is set to zero:
\[
\chi_a=0
\quad\Longrightarrow\quad
\frac{\chi_a}{s_a}=0.
\]
The supported exact inverse is
\[
G_\tau^+
=
V\diag\left(\frac{\chi_a}{s_a}\right)V^\top,
\]
with the convention \(\chi_a/s_a=0\) when \(\chi_a=0\). In components,
\[
(G_\tau^+)_{ij}
=
\sum_{a=1}^{n}V_{ia}\frac{\chi_a}{s_a}V_{ja}.
\]

\phys{A retained eigenmode is a tangent-coordinate combination that changes the
state enough to be trusted. A truncated mode is treated as geometrically
unresolved, so the inverse does not amplify that direction.}

\geom{The cutoff draws a floor under the metric ellipsoid. Axes above the floor
remain available for inversion; axes below the floor are removed from the
inverse, so the solve does not chase directions the ansatz barely resolves.}

\micro{The damped implementation inverse changes the denominator of retained
modes.}

\deriv
Let \(\epsilon\ge0\). The damped supported inverse is
\[
G_{\tau,\epsilon}^{-1}
=
V\diag\left(\frac{\chi_a}{s_a+\epsilon}\right)V^\top.
\]
Thus
\[
\dot\theta
=G_{\tau,\epsilon}^{-1}f
=
\sum_{a=1}^{n}
V_{\cdot a}
\frac{\chi_a}{s_a+\epsilon}
(V_{\cdot a}^\top f).
\]

\phys{Damping makes each retained inverse denominator larger. This reduces the
amplification of small retained modes and can reduce unstable parameter
velocity.}

\micro{Ridge changes the matrix before the eigenmode solve.}

\deriv
Let \(\Lambda\succeq0\). The propagated solve uses
\[
\begin{aligned}
G_\Lambda&=G+\Lambda,\\
\dot\theta_\ell
&=(G+\Lambda_\ell)_{\tau_\ell,\epsilon_\ell}^{-1}f.
\end{aligned}
\]
If \(\Lambda_\ell=\lambda_\ell\Id\), then
\[
v^\top(G+\lambda_\ell\Id)v
=v^\top Gv+\lambda_\ell\norm v^2.
\]

\phys{Ridge penalizes large parameter velocity in coordinate space. It is a
numerical stabilizer for the solve, not an extra Hamiltonian force.}

\micro{The supported inverse and the state-space trust region control different
objects in the local Taylor model.}

\deriv
For a small coordinate displacement \(\delta\theta\),
\[
\ket{\psi(\theta+\delta\theta)}
=
\ket{\psi(\theta)}+T\delta\theta+O(\norm{\delta\theta}^2).
\]
The supported inverse acts on metric eigenmodes:
\[
\begin{aligned}
s_a\ge\tau s_{\max}
&\Longrightarrow
v_a\ \text{retained},\\
s_a<\tau s_{\max}
&\Longrightarrow
v_a\ \text{truncated}.
\end{aligned}
\]
The state-space trust condition acts after a velocity has been chosen:
\[
\delta\theta=\Delta t\,\dot\theta,
\qquad
\Delta t\,\norm{T\dot\theta}\le\tau_\Delta.
\]

\phys{The inverse decides which coordinate axes are identifiable enough to
invert. The trust condition decides whether the accepted velocity moves the
state too far for the local linear model. A well-conditioned inverse can still
require a smaller time step.}

\section*{8. Realized drift and numerical miss}

\micro{Numerical miss measures how much explained tangent motion is lost because
of the chosen stabilized inverse.}

\deriv
Let
\[
\dot\theta_\ell=(G+\Lambda_\ell)_{\tau_\ell,\epsilon_\ell}^{-1}f.
\]
The residual produced by this specific velocity is
\[
\begin{aligned}
m(\dot\theta_\ell)
&=\norm{T\dot\theta_\ell-\ket b}^2\\
&=\dot\theta_\ell^\top G\dot\theta_\ell
-2f^\top\dot\theta_\ell+\norm{\ket b}^2.
\end{aligned}
\]
The realized explained drift is
\[
\begin{aligned}
E_\ell
&=\norm{\ket b}^2-m(\dot\theta_\ell)\\
&=2f^\top\dot\theta_\ell
-\dot\theta_\ell^\top G\dot\theta_\ell.
\end{aligned}
\]
The best supported explained drift is
\[
E_\star=f^\top G^+f.
\]
The numerical miss is
\[
\boxed{
\rho_{\rm num}(\ell)
=
\frac{[E_\star-E_\ell]_+}{\norm{\ket b}^2}.
}
\]

\phys{If \(\rho_{\rm num}\) is large, the tangent span may contain useful
directions but the current inverse convention is not realizing them. This is a
solve problem, not automatically an append problem.}

\micro{The realized drift uses \(G\), not \(G+\Lambda_\ell\), because the
question is physical tangent motion, not self-certification by the regularizer.}

\deriv
\[
\begin{aligned}
E_\ell
&=2f^\top\dot\theta_\ell-\dot\theta_\ell^\top G\dot\theta_\ell,\\
\dot\theta_\ell
&=(G+\Lambda_\ell)_{\tau_\ell,\epsilon_\ell}^{-1}f.
\end{aligned}
\]
If one used \(G+\Lambda_\ell\) in the realized term, then
\[
2f^\top\dot\theta_\ell
-\dot\theta_\ell^\top(G+\Lambda_\ell)\dot\theta_\ell
\]
would include the penalty that chose \(\dot\theta_\ell\), rather than only the
state-space velocity \(T\dot\theta_\ell\).

\phys{The propagated state only sees \(T\dot\theta\). The ridge penalty helps
choose the velocity, but the physical miss must be evaluated through \(T\) and
therefore through \(G=\Real(T^\dagger T)\).}

\section*{9. State-space step motion}

\micro{A stable inverse can still move the state too far for the selected time
step.}

\deriv
Let
\[
\ket{u_\ell}=T\dot\theta_\ell\in\Hh,
\qquad
\Delta t_\ell>0.
\]
First-order propagation gives
\[
\ket{\psi(t+\Delta t_\ell)}
=\ket{\psi(t)}+\Delta t_\ell\ket{u_\ell}+O(\Delta t_\ell^2).
\]
Thus
\[
\norm{\ket{\psi(t+\Delta t_\ell)}-\ket{\psi(t)}}
=
\Delta t_\ell\norm{\ket{u_\ell}}+
O(\Delta t_\ell^2).
\]
The state-motion guard is
\[
\Delta t_\ell\norm{T\dot\theta_\ell}\le\tau_\Delta.
\]

\phys{This guard checks the physical size of the proposed state movement, not
the size of raw parameter movement. If it fails, local subdivision is a natural
cure.}

\micro{Local subdivision reduces the first-order state displacement per
substep.}

\deriv
For \(r\) equal substeps,
\[
\delta t=\Delta t/r.
\]
The one-substep first-order motion is
\[
\delta t\,\norm{\ket{u_\ell}}
=\frac{\Delta t}{r}\norm{\ket{u_\ell}}.
\]
To satisfy \(\delta t\norm{\ket{u_\ell}}\le\tau_\Delta\), it is enough that
\[
r\ge \frac{\Delta t\norm{\ket{u_\ell}}}{\tau_\Delta}.
\]
The Paper-II scheduler uses a severity-derived dyadic response:
\[
r=2^q.
\]

\phys{Subdivision is not artificial damping. It keeps the same local vector
field idea but asks the integrator to take smaller local steps through a
difficult interval.}

\section*{10. Temporal state-space kink}

\micro{The kink diagnostic compares represented state velocities at adjacent
accepted checkpoints.}

\deriv
Let
\[
\ket{u_k}=T_k\dot\theta_k,\qquad
\ket{u_{k-1}^{\rm acc}}=T_{k-1}\dot\theta_{k-1}^{\rm acc}.
\]
Then
\[
\eta_k
=
\frac{\norm{\ket{u_k}-\ket{u_{k-1}^{\rm acc}}}^2}
{\tfrac12(\norm{\ket{b_k}}^2+\norm{\ket{b_{k-1}}}^2)}.
\]
The guard is
\[
\eta_k\le\tau_{\rm kink}.
\]

\phys{A kink is a jump in the physical tangent velocity. It is evaluated in
state space through \(T\dot\theta\), not by comparing coordinate vectors.}

\section*{11. Repair entry}

\micro{Repair opens when the base numerical convention is unhealthy in
state-space terms.}

\deriv
A rung is one numerical convention for the same local McLachlan problem:
\[
\ell=(\Lambda_\ell,\tau_\ell,\epsilon_\ell,\Delta t_\ell).
\]
Here \(\Lambda_\ell\) is the ridge matrix, \(\tau_\ell\) is the retained-mode
threshold, \(\epsilon_\ell\) is the solve damping, and \(\Delta t_\ell\) is the
time step for that local proposal. The base rung \(\ell_0\) is the default
convention before repair. At the base rung, define
\[
\begin{aligned}
g_\rho&=\mathbf 1[\rho_{\rm num}(\ell_0)>\tau_{\rm num}],\\
g_\Delta&=\mathbf 1[\Delta t_{\ell_0}\norm{T\dot\theta_{\ell_0}}>\tau_\Delta],\\
g_t&=\mathbf 1[\eta_k(\ell_0)>\tau_{\rm kink}].
\end{aligned}
\]
The symbol \(\mathbf 1[\cdot]\) is an indicator: it equals \(1\) when the
condition inside brackets is true and \(0\) otherwise. Let
\(g_{\rm inv}\in\{0,1\}\) mark invalid numerical data. Here invalid means the
needed arrays do not define a real finite solve: NaN/Inf values, missing
objects, nonsymmetric real Gram data after preprocessing, or incompatible
matrix/vector dimensions. Repair entry is checked only after the checkpoint
data are finite and dimensionally compatible.
Then
\[
g_{\rm inv}=0.
\]
Repair entry is
\[
g_{\rm inv}\vee g_\rho\vee g_\Delta\vee g_t.
\]

\phys{The three non-invalid lanes are: numerical miss, state-motion size, and
temporal kink. The condition number modifies candidate risk; it is not a
global repair trigger in the noiseless diagnostic policy.}

\micro{The severity scalar determines how broad the repair response should be.}

\deriv
The severity scalar \(m\) is a dimensionless maximum over normalized defects:
\(m\le1\) means no listed defect exceeds its threshold, while \(m>1\) measures
how far the worst defect is beyond threshold.
\[
\begin{aligned}
m
=
\max\Bigg\{
&
\frac{\rho_{\rm num}(\ell_0)}{\tau_{\rm num}},
\\
&
\frac{\Delta t_{\ell_0}\norm{T\dot\theta_{\ell_0}}}{\tau_\Delta},
\\
&
\sqrt{\frac{\eta_k(\ell_0)}{\tau_{\rm kink}}}
\Bigg\}.
\end{aligned}
\]
The scheduler breadth is
\[
q=\left\lceil\log_2(m)\right\rceil_+,
\qquad
\lceil x\rceil_+=\max\{0,\lceil x\rceil\}.
\]

\phys{A larger normalized defect requests stronger local subdivision, broader
inverse search, or longer hold behavior. The square root on \(\eta_k\)
puts the kink's squared norm back on a velocity-size scale.}

\section*{12. Candidate repair set}

\micro{Paper II does not use an ordered first-passing repair list. It evaluates
a finite set of candidates.}

\deriv
For repair, a candidate is a proposed numerical rung for the same current
state, tangent matrix, and drift vector. Let
\[
\ell=(\Lambda_\ell,\tau_\ell,\epsilon_\ell,\Delta t_\ell).
\]
Let \(C\) be the finite candidate set considered at the checkpoint:
\[
C=\{\ell_0\}\cup C_{\rm sub}\cup C_{\rm inv}.
\]
Here \(\{\ell_0\}\) keeps the base convention available,
\(C_{\rm sub}\) contains time-subdivision candidates, and \(C_{\rm inv}\)
contains inverse-convention candidates. The subdivision candidates are present
when \(g_\Delta=1\) or \(g_t=1\):
\[
C_{\rm sub}
=
\{\ell_{\Delta t/2},\ell_{\Delta t/4},\ldots,\ell_{\Delta t/2^q}\}.
\]
The notation \(\ell_{\Delta t/2^j}\) means the same inverse settings as the
base rung, but with the local step divided by \(2^j\).
The inverse candidates are present when \(g_\rho=1\):
\[
C_{\rm inv}
=
C_\uparrow\cup C_\downarrow.
\]
Here \(C_\uparrow\) increases regularization/truncation/damping, while
\(C_\downarrow\) safely relaxes them.

\phys{The same numerical miss can come from under-regularization or
over-regularization. The candidate set tests nearby directions instead of
assuming one cure direction.}

\micro{Acceptability is a state-space condition on each candidate rung.}

\deriv
An acceptable rung is a candidate whose realized state-space diagnostics all
satisfy their declared thresholds.
For \(\ell\in C\),
\[
\begin{aligned}
\dot\theta_\ell
&=(G+\Lambda_\ell)_{\tau_\ell,\epsilon_\ell}^{-1}f,\\
\ket{u_\ell}&=T\dot\theta_\ell.
\end{aligned}
\]
The finite acceptable set is
\[
\begin{aligned}
A=\{\ell\in C:\;&
\rho_{\rm num}(\ell)\le\tau_{\rm num},\\
&
\Delta t_\ell\norm{\ket{u_\ell}}\le\tau_\Delta,\\
&
\eta_k(\ell)\le\tau_{\rm kink},\\
&
\ell\text{ is finite}\}.
\end{aligned}
\]

\phys{A candidate rung has to actually improve the realized state-space
behavior. Merely being more damped or more aggressive does not make it better.}

\section*{13. Repair selection}

\micro{The accepted repair is the least invasive acceptable candidate.}

\deriv
Least invasive means lowest scalar score \(J(\ell)\) among acceptable rungs.
The score combines a declared intervention cost with normalized defects. The
weights \(w_c,w_\rho,w_\Delta,w_t,w_\eta,w_\kappa\ge0\) are runtime settings,
not derived physics quantities.
For \(\ell\in C\), let
\[
\begin{aligned}
J(\ell)
={}&
w_c C_\ell
+w_\rho\frac{\rho_{\rm num}(\ell)}{\tau_{\rm num}}
+w_\Delta
\frac{\Delta t_\ell\norm{\ket{u_\ell}}}{\tau_\Delta}\\
&+
w_t\frac{\eta_k(\ell)}{\tau_{\rm kink}}
+w_\eta\frac{\eta_{\rm cand}(\ell)}{\tau_{\rm kink}}
+w_\kappa\phi_\kappa(\ell).
\end{aligned}
\]
Here \(C_\ell\) is the intervention cost of candidate \(\ell\): small for
staying near the base convention and larger for stronger damping, truncation
changes, or subdivision. The function \(\phi_\kappa(\ell)\ge0\) is a
conditioning-risk penalty computed from the retained condition number of the
candidate solve.
The selected rung is
\[
\ell^\star\in\argmin_{\ell\in A}J(\ell).
\]
If \(A=\varnothing\) in a diagnostic run,
\[
\ell^\star\in\argmin_{\ell\in C_{\rm finite}}J(\ell),
\]
and the checkpoint is marked unsupported.
The set \(C_{\rm finite}\subseteq C\) contains candidates whose arrays and
velocities are finite even though one or more diagnostic thresholds failed.

\phys{Unsupported means the run produced diagnostic data beyond a numerical
health threshold. It does not mean the trajectory had a runtime error.}

\micro{The candidate-kink term compares a candidate against the base solve at
the same checkpoint.}

\deriv
This term is local in time: it compares two candidate velocities at checkpoint
\(k\), not velocities at two different checkpoints.
Let
\[
\ket{u_\ell}=T\dot\theta_\ell,\qquad \ket{u_0}=T\dot\theta_0.
\]
Then
\[
\eta_{\rm cand}(\ell)
=
\frac{\norm{\ket{u_\ell}-\ket{u_0}}^2}{\norm{\ket b}^2}.
\]

\phys{This term discourages a repair that creates a large same-checkpoint
velocity jump relative to the base convention, unless the other terms justify
it.}

\section*{14. Hold and release}

\micro{After a repair is accepted, the method should not immediately flip back
to the base inverse if that creates another velocity jump.}

\deriv
A held convention is an accepted non-base rung reused at later checkpoints
until release is safe. Let the held repaired convention be \(\ell^\star\). At a
later checkpoint,
\[
\begin{aligned}
\ket{u_{\rm hold}}&=T\dot\theta_{\ell^\star},\\
\ket{u_{\rm base}}&=T\dot\theta_0.
\end{aligned}
\]
The return-to-base comparison is
\[
\eta_{\rm rel}
=
\frac{\norm{\ket{u_{\rm hold}}-\ket{u_{\rm base}}}^2}{\norm{\ket b}^2}.
\]
Release requires
\[
\eta_{\rm rel}\le \tau_{\rm rel},
\qquad
\tau_{\rm rel}=\alpha_{\rm rel}\tau_{\rm kink},
\qquad
0<\alpha_{\rm rel}<1.
\]

\phys{The release threshold is stricter than the original kink threshold. This
prevents a repaired inverse from turning on for one checkpoint and then turning
off in a way that creates rough state-space velocity.}

\micro{The patience count grows with the size of the accepted kink.}

\deriv
The patience count \(p\) is the number of future checkpoints for which the
method may keep the repaired convention before attempting release. Let
\[
\eta_{\rm acc}
=
\frac{\norm{\ket{u^+}-\ket{u^-}}^2}{\norm{\ket b}^2},
\]
where \(\ket{u^-}\) and \(\ket{u^+}\) are the pre-repair and accepted post-repair
represented velocities. Then
\[
\begin{aligned}
p
=
\Bigg\lceil
p_{\min}
&+(p_{\max}-p_{\min})\\
&\times
\min\left\{
1,
\frac{[\sqrt{\eta_{\rm acc}}-\sqrt{\tau_{\rm kink}}]_+}
{\sqrt{\tau_{\rm kink}}\chi_{\rm kink}}
\right\}
\Bigg\rceil.
\end{aligned}
\]

\phys{The square roots again put the kink severity on a velocity scale. A severe
repair is held longer before the base convention can return.}

\section*{15. Append: zero state change, new tangent directions}

\micro{Append adds directions to the tangent space without changing the state at
the instant of insertion.}

\deriv
An append block \(R\) is a set of proposed new generators inserted with zero
initial amplitude. If the block has \(r\) new real coordinates, write
\(\phi\in\R^r\). The phrase zero amplitude means \(\phi=0\) at insertion, so
the state is unchanged before the tangent space is re-evaluated. Let
\[
U_R(\phi):\Hh\to\Hh,\qquad U_R(0)=\Id.
\]
The augmented state is
\[
\ket{\psi_R(\theta,\phi)}=U_R(\phi)\ket{\psi(\theta)}.
\]
At zero amplitude,
\[
\ket{\psi_R(\theta,0)}=\ket{\psi(\theta)}.
\]
The candidate tangent columns are
\[
\ket{u_a}=\Pi_\psi\,\partial_{\phi_a}\ket{\psi_R(\theta,\phi)}\big|_{\phi=0},
\qquad
1\le a\le r.
\]
Collect them in
\[
U=
\begin{bmatrix}
| & & |\\
\ket{u_1}&\cdots&\ket{u_r}\\
| & & |
\end{bmatrix}
\in\C^{d\times r}.
\]

\phys{Append does not jump the state. It changes which instantaneous velocities
the McLachlan solve can represent after the checkpoint.}

\micro{For a Hermitian appended generator \(A\), the zero-angle tangent has the
same physical form as a McLachlan direction.}

\deriv
Let
\[
U_R(\phi)=\exp(-\ii\phi A),
\qquad A=A^\dagger.
\]
Then
\[
\begin{aligned}
\partial_\phi U_R(\phi)
&=-\ii A\exp(-\ii\phi A),\\
\partial_\phi\ket{\psi_R(\theta,\phi)}\big|_{\phi=0}
&=-\ii A\ket{\psi},\\
\ket u&=\Pi_\psi(-\ii A\ket{\psi}).
\end{aligned}
\]

\phys{A drive-aligned density augmentation uses this exact logic: it adds the
tangent \(-\ii D\ket\psi\) for the drive operator \(D\), while leaving \(H(t)\) and
the state unchanged at zero angle.}

\section*{16. Augmented least squares}

\micro{Append gain is a before--after least-squares comparison on the augmented
tangent matrix.}

\deriv
Let
\[
T_R=[T~U]\in\C^{d\times(n+r)}.
\]
The augmented coordinate vector \(w\) stacks old McLachlan velocity coordinates
and candidate-block velocity coordinates. Let
\[
w=
\begin{bmatrix}
\dot\theta\\
\alpha
\end{bmatrix}
\in\R^{n+r},
\qquad
\dot\theta\in\R^n,\quad \alpha\in\R^r.
\]
The augmented represented velocity is
\[
T_Rw=T\dot\theta+U\alpha.
\]
The augmented objective is
\[
m_R(w)=\norm{T_Rw-\ket b}^2.
\]
Define
\[
G_R=\Real(T_R^\dagger T_R),
\qquad
f_R=\Real(T_R^\dagger\ket b).
\]
Then
\[
m_R(w)=w^\top G_Rw-2w^\top f_R+\norm{\ket b}^2.
\]

\phys{The old coordinates and new candidate coordinates are fitted together.
The append score is not the raw overlap of \(U\) with \(\ket b\); it is the improvement
after the whole augmented tangent space refits.}

\micro{The augmented matrices have a block structure that exposes old--new
coupling.}

\deriv
\[
G_R
=
\begin{bmatrix}
G & B\\
B^\top & C
\end{bmatrix},
\qquad
f_R=
\begin{bmatrix}
f\\
q
\end{bmatrix},
\]
where
\[
\begin{aligned}
G&=\Real(T^\dagger T)\in\R^{n\times n},\\
B&=\Real(T^\dagger U)\in\R^{n\times r},\\
C&=\Real(U^\dagger U)\in\R^{r\times r},\\
f&=\Real(T^\dagger\ket b)\in\R^n,\\
q&=\Real(U^\dagger\ket b)\in\R^r.
\end{aligned}
\]
In components,
\[
\begin{aligned}
B_{ia}
&=\Real\sum_{\mu=1}^{d}T_{\mu i}^*U_{\mu a},\\
C_{ab}
&=\Real\sum_{\mu=1}^{d}U_{\mu a}^*U_{\mu b},\\
q_a
&=\Real\sum_{\mu=1}^{d}U_{\mu a}^*b_\mu.
\end{aligned}
\]

\phys{\(B\) measures how much the candidate tangent block overlaps the old
tangent directions. \(C\) is the candidate internal Gram matrix. \(q\) is the
candidate's direct force vector.}

\section*{17. Append gain}

\micro{The append gain is the increase in explained drift after solving the
augmented normal equation.}

\deriv
The notation \(G_{\rm stab}^{-1}\) means the active stabilized inverse: a
supported inverse with the current retained-mode threshold, ridge, and damping.
For this block it is shorthand, not an ordinary matrix inverse.
Old explained drift:
\[
E_{\rm old}=f^\top G_{\rm stab}^{-1}f.
\]
Augmented explained drift:
\[
E_R=f_R^\top G_{R,{\rm stab}}^{-1}f_R.
\]
Append gain:
\[
\rho_{\rm gain}(R)
=
\frac{[E_R-E_{\rm old}]_+}{\norm{\ket b}^2}.
\]
Expanded:
\[
\rho_{\rm gain}(R)
=
\frac{
\left[
\begin{bmatrix}f\\q\end{bmatrix}^{\!\top}
\begin{bmatrix}G&B\\B^\top&C\end{bmatrix}_{\rm stab}^{-1}
\begin{bmatrix}f\\q\end{bmatrix}
-f^\top G_{\rm stab}^{-1}f
\right]_+
}{\norm{\ket b}^2}.
\]

\phys{This is Paper II's append score in simplified notation. A singleton is the
case \(r=1\). A batch is the case \(r>1\).}

\micro{The post-append structural miss is the same projection residual on the
augmented tangent span.}

\deriv
\[
\rho(T_R,\ket b)
=
\frac{
\left[
\norm{\ket b}^2
-f_R^\top G_R^+f_R
\right]_+
}{\norm{\ket b}^2}.
\]
When \(R=\varnothing\),
\[
T_R=T,\qquad
G_R=G,\qquad
f_R=f.
\]
Thus
\[
\rho(T,\ket b)
=
\frac{
\left[
\norm{\ket b}^2-f^\top G^+f
\right]_+
}{\norm{\ket b}^2}.
\]

\phys{The empty-set value is the current structural miss. Candidate append
blocks are tested by how much they reduce that miss and how much stabilized
gain they provide.}

\section*{18. Schur novelty and eliminated old coordinates}

\micro{The Schur expression appears when the old coordinates are optimized out
and only the candidate's leftover direction remains.}

\deriv
The theorem used here is the Schur-complement block-elimination identity. In
this context, eliminating a coordinate means solving for that coordinate's
least-squares optimum and substituting it back into the remaining equations. If a
block matrix has an invertible old-support block \(G\), then eliminating the old
variables replaces the candidate block \(C\) by
\[
S=C-B^\top G^{-1}B.
\]
Assume \(G\) is invertible for the derivation. The augmented normal equations
are
\[
\begin{bmatrix}
G&B\\
B^\top&C
\end{bmatrix}
\begin{bmatrix}
\dot\theta\\
\alpha
\end{bmatrix}
=
\begin{bmatrix}
f\\
q
\end{bmatrix}.
\]
Equivalently,
\[
\begin{aligned}
G\dot\theta+B\alpha&=f,\\
B^\top\dot\theta+C\alpha&=q.
\end{aligned}
\]
Solve the first line for \(\dot\theta\):
\[
\dot\theta=G^{-1}f-G^{-1}B\alpha.
\]
Substitute into the second line:
\[
\begin{aligned}
B^\top(G^{-1}f-G^{-1}B\alpha)+C\alpha&=q,\\
(C-B^\top G^{-1}B)\alpha&=q-B^\top G^{-1}f.
\end{aligned}
\]
Define
\[
S=C-B^\top G^{-1}B,
\qquad
w=q-B^\top G^{-1}f.
\]
Then
\[
S\alpha=w.
\]

\phys{The old tangent span is allowed to explain both the drift and the
candidate. After eliminating the old coefficients, \(S\) is the candidate's
leftover Gram geometry and \(w\) is its leftover force.}

\geom{Picture the old tangent columns as a plane and the candidate block as new
arrows arriving near that plane. Schur first lets the old plane absorb every
shadow it can. The remaining arrows are the candidate's leftover directions:
\(S\) measures their geometry, and \(w\) measures the drift shadow landing on
them.}

\micro{The Schur gain is the explained drift inside that leftover candidate
space.}

\deriv
With \(\alpha=S^{-1}w\),
\[
\begin{aligned}
\Delta_{\rm Schur}
&=w^\top S^{-1}w\\
&=\alpha^\top S\alpha.
\end{aligned}
\]
Using the block elimination above,
\[
E_R-E_{\rm old}
=w^\top S^{-1}w.
\]
Thus
\[
\rho_{\rm gain}(R)
=
\frac{[w^\top S^{-1}w]_+}{\norm{\ket b}^2}
\]
when the Schur shortcut uses the same retained-support and regularization
convention as the direct grouped solve.

\phys{Schur is not a separate physics objective. It is a compact way to compute
the same grouped before--after least-squares gain after old coordinates have
been optimized out.}

\caution{If whitening, compression, retained-mode truncation, or ridge differs
between the direct grouped solve and Schur shortcut, the direct grouped gain is
the authority. The Schur value can still screen or confirm candidates.}

\section*{19. Cost-weighted append and the ladder}

\micro{Cost weighting divides gain by a hardware-cost factor; it does not turn
a weak tangent candidate into an admissible one.}

\deriv
Let \(K_{\rm hw}(R)\ge1\) be the hardware-cost factor of append block \(R\),
such as a compiled two-qubit-count or estimator-work proxy. Let
\(\alpha_{\rm app}\ge0\) be the exponent controlling how strongly cost reduces
the append score. The cost-weighted append score is
\[
\Xi_{\rm app}(R)
=
\frac{\rho_{\rm gain}(R)}
{K_{\rm hw}(R)^{\alpha_{\rm app}}}.
\]
The direct-admission test has the form
\[
\begin{aligned}
\rho(T,\ket b)&>\tau_{\rm miss},\\
\rho_{\rm gain}(R)&\ge\tau_{\rm gain},\\
\Xi_{\rm app}(R)&\ge\tau_{\rm app},\\
S_R&\succeq_{\rm retained}\tau_{\rm Schur},\\
G_R^{-1}f_R&\text{ finite under the active solve convention.}
\end{aligned}
\]
The notation \(S_R\succeq_{\rm retained}\tau_{\rm Schur}\) means that the
candidate's Schur-reduced Gram matrix has enough retained metric support: after
the same retained-mode convention is applied, the accepted candidate directions
are not numerically null below the Schur threshold.

\phys{The candidate must be useful, affordable, geometrically retained, and
numerically sane before commit. Exact reference trajectories are not part of
this decision.}

\micro{The combinatorial batch ladder scores joint sets at each cardinality.}

\deriv
Let \(P\) be the available candidate pool: the finite set of append blocks that
can be tested at the checkpoint. A ladder rung \(q\) is a candidate batch size.
At rung \(q\),
\[
\mathcal R_q=\{R\subseteq P:\ |R|=q\}.
\]
The selected set at rung \(q\) is
\[
R_q^\star\in\argmax_{R\in\mathcal R_q}\Xi_{\rm app}(R)
\]
subject to the confirmation checks. The ladder escalates only while the miss
remains active and smaller accepted sets did not satisfy the declared
admission rule.

\phys{A size-five batch means ``at most five, tested through the ladder''
when the implementation follows the Paper-II ladder. It should not mean forcing
exactly five candidates regardless of miss and confirmation.}

\section*{20. Failed append search reuse}

\micro{A failed append search should not force repeated full rescoring until
the local model has changed enough.}

\deriv
At a failed search checkpoint \(k_0\), store a certificate. A certificate is the
data needed to decide whether the old failed-search conclusion is still
applicable:
\[
\mathcal C_{k_0}
=
(J_{k_0},P_{k_0},\pi_{k_0},r_{k_0},G_{k_0},f_{k_0},U^\star_{k_0}).
\]
Here \(J_{k_0}\) is the active coordinate support, \(P_{k_0}\) is the candidate
pool, \(\pi_{k_0}\) is the parameterization policy, \(r_{k_0}\) is the maximum
batch rung tested, and \(U^\star_{k_0}\) is the best candidate block found
before failure. The support symbol \(J_{k_0}\) is unrelated to the repair score
\(J(\ell)\) in Section 13.
At a later checkpoint \(k\), immediate invalidation occurs if
\[
J_k\ne J_{k_0}
\quad\text{or}\quad
P_k\ne P_{k_0}
\quad\text{or}\quad
\pi_k\ne\pi_{k_0}
\quad\text{or}\quad
r_k\ne r_{k_0}.
\]
If not invalidated, use local drift metrics such as
\[
\begin{aligned}
d_G(k,k_0)
&=\norm{\widetilde G_{k_0}^{-1/2}
(G_k-G_{k_0})
\widetilde G_{k_0}^{-1/2}},\\
d_f(k,k_0)
&=
\frac{\norm{\widetilde G_{k_0}^{-1/2}(f_k-f_{k_0})}}
{1+\norm{\widetilde G_{k_0}^{-1/2}f_{k_0}}}.
\end{aligned}
\]
The matrix \(\widetilde G_{k_0}\) is a stabilized reference metric used to make
changes in \(G\) and \(f\) dimensionless before comparing checkpoints.
The search reopens when the combined drift exceeds the declared retry rule.

\phys{This is not a total append cap. It is a local model-staleness rule:
skip full candidate rescoring only while the old failed-search conclusion is
still geometrically credible.}

\section*{21. Prune as the reverse least-squares comparison}

\micro{Prune asks how much explained drift is lost after deleting directions
and refitting the survivors.}

\deriv
A deletion batch \(D\) is a proposed set of active coordinates to remove
together. Refit means the retained coordinates are re-solved after deletion;
the deleted block is not simply zeroed while old velocities are held fixed.
Partition the active tangent matrix as
\[
T=[T_D~T_R],
\]
where \(T_D\in\C^{d\times r}\) are directions proposed for deletion and
\(T_R\in\C^{d\times(n-r)}\) are the retained directions. The full primitives
are
\[
G=\Real(T^\dagger T),
\qquad
f=\Real(T^\dagger\ket b).
\]
The retained-only primitives are
\[
G_R=\Real(T_R^\dagger T_R),
\qquad
f_R=\Real(T_R^\dagger\ket b).
\]
Full explained drift:
\[
E_{\rm full}=f^\top G_{\rm stab}^{-1}f.
\]
Retained explained drift:
\[
E_R=f_R^\top G_{R,{\rm stab}}^{-1}f_R.
\]
Deletion loss:
\[
\boxed{
L(D)
=
\frac{[E_{\rm full}-E_R]_+}{\norm{\ket b}^2}.
}
\]

\phys{Small deletion loss means the remaining tangent directions can still
explain nearly the same instantaneous drift after the deleted block is removed
and the survivors refit.}

\micro{The singleton prune case is the one-block version of the same grouped
loss.}

\deriv
If \(D=\{d\}\), then
\[
T=[t_d~T_R],
\]
and
\[
L(d)
=
\frac{
\left[
f^\top G_{\rm stab}^{-1}f
-f_R^\top G_{R,{\rm stab}}^{-1}f_R
\right]_+
}{\norm{\ket b}^2}.
\]
If \(D=\{d_1,\ldots,d_q\}\), the same formula holds with all deleted columns
removed at once.

\phys{Batched prune is not the sum of singleton prune losses. It is the grouped
loss after removing the whole proposed set and refitting the retained support.}

\section*{22. Prune Schur form}

\micro{The prune Schur expression is the same elimination geometry with the sign
reversed.}

\deriv
This again uses the Schur-complement block-elimination identity, now after
splitting the support into deleted coordinates \(D\) and retained coordinates
\(R\).
Partition the full Gram matrix and force as
\[
G=
\begin{bmatrix}
G_{DD}&G_{DR}\\
G_{RD}&G_{RR}
\end{bmatrix},
\qquad
f=
\begin{bmatrix}
f_D\\
f_R
\end{bmatrix}.
\]
Eliminate retained coordinates first. The deletion-side Schur objects are the
leftover Gram matrix and leftover force for the deleted block after the retained
block has made its best local fit:
\[
\begin{aligned}
S_D^-&=G_{DD}-G_{DR}G_{RR}^{-1}G_{RD},\\
w_D^-&=f_D-G_{DR}G_{RR}^{-1}f_R.
\end{aligned}
\]
Under matched inverse conventions,
\[
L(D)
=
\frac{[(w_D^-)^\top(S_D^-)^{-1}w_D^-]_+}{\norm{\ket b}^2}.
\]

\geom{For prune, reverse the picture. The retained arrows form the plane that
stays. The proposed deleted block is safe only when its leftover direction,
after the retained plane explains what it can, is small and carries little
force shadow.}

\phys{If the deleted block is already explained by the retained block, its
Schur leftover is small and its loss is small. This is a local tangent-plane
statement, not a full future-trajectory certificate.}

\section*{23. Why prune needs history}

\micro{A direction can be redundant at one checkpoint and useful later, so prune
needs persistence.}

\deriv
The history set \(H_D\) is the finite collection of recent checkpoints used to
test whether deletion batch \(D\) has been persistently low-loss, not only
accidentally low-loss at the present checkpoint. For a deletion batch \(D\),
define checkpoint losses
\[
L_\ell(D)
=
\frac{[E_{\rm full,\ell}-E_{R,\ell}]_+}{\norm{\ket{b_\ell}}^2}.
\]
For weights \(w_\ell\ge0\),
\[
\sum_{\ell\in H_D}w_\ell=1,
\]
the history loss is
\[
\overline L_{\rm hist}(D)
=
\sum_{\ell\in H_D}w_\ell L_\ell(D).
\]
The parameter \(\lambda_{\rm hist}\ge0\) controls how strongly historical loss
enters the score, and \(\epsilon_L>0\) prevents division by zero when both
current and historical losses are tiny. The cost-weighted prune score can be
written
\[
\Xi_{\rm prune}(D)
=
\frac{K_{\rm hw}(D)^{\alpha_{\rm prune}}}
{L(D)+\lambda_{\rm hist}\overline L_{\rm hist}(D)+\epsilon_L}.
\]

\phys{Current low loss nominates a deletion. Historical low loss makes the
nomination more credible. Projection protects the current state, a short-horizon
no-harm check tests a few immediate future steps, cooldown prevents repeated
rapid edits of the same support, and rollback means restoring the previous
support if the reduced run fails its safety checks.}

\section*{24. Reduced-state projection after deletion}

\micro{Accepted prune cannot be unverified deletion; the reduced ansatz must be
projected back near the incumbent state.}

\deriv
Let \(\theta\in\R^n\) be the incumbent coordinate vector. Let \(D\) be deleted,
and let the reduced coordinate vector be \(\xi\in\R^{n-r}\). The reduced state is
\[
\ket{\psi_{-D}(\xi)}=U_{-D}(\xi)\ket{\psi_{\rm ref}}.
\]
The projection objective has the form
\[
\xi^\star\in\argmin_{\xi\in\R^{n-r}}
\Phi_D(\xi),
\]
with
\[
\begin{aligned}
\Phi_D(\xi)
={}&
d_{\rm FS}^2(\psi_{-D}(\xi),\psi(\theta))\\
&+\lambda_Y\norm{Y(\xi)-Y(\theta)}_{\Sigma_Y^{-1}}^2\\
&+\lambda_u\norm{\ket{u(\xi)}-\ket{u(\theta)}}^2.
\end{aligned}
\]
Here \(d_{\rm FS}\) is Fubini--Study distance, the projective Hilbert-space
distance between normalized states after global phase is removed. The vector
\(Y(\xi)\) collects tracked observables, \(\Sigma_Y\) is their positive
weighting/covariance matrix, and
\(\norm{x}_{\Sigma_Y^{-1}}^2=x^\top\Sigma_Y^{-1}x\). The constants
\(\lambda_Y,\lambda_u\ge0\) weight observable preservation and represented
velocity preservation.
The committed reduced state is
\[
\ket{\psi_{-D}(\xi^\star)}.
\]

\phys{Deletion changes the ansatz support. The projection step chooses reduced
coordinates that preserve the state, tracked observables, and represented
velocity as well as the reduced support allows.}

\section*{25. Patch after pure append and pure prune}

\micro{The support patch is simple after pure append and pure prune are already
understood.}

\deriv
A support patch is one checkpoint edit to the active coordinate support. It has
a deletion side \(D\) and an append side \(R\), written
\(\nu=(D,R)\). The patched tangent matrix is
\[
T_\nu=[T_{-D}~U_R],
\qquad
\nu=(D,R).
\]
Here \(T_{-D}\) means the old tangent matrix after removing the columns indexed
by \(D\), and \(U_R\) is the zero-amplitude tangent block for the appended set
\(R\).
Define
\[
G_\nu=\Real(T_\nu^\dagger T_\nu),
\qquad
f_\nu=\Real(T_\nu^\dagger\ket b).
\]
The patch explained-drift change is
\[
\Delta_{\rm patch}(\nu)
=
\frac{
f_\nu^\top G_{\nu,{\rm stab}}^{-1}f_\nu
-f^\top G_{\rm stab}^{-1}f
}{\norm{\ket b}^2}.
\]
The limiting cases are
\[
\begin{array}{ccl}
D=\varnothing,\ R=\varnothing &:& \text{stay},\\
D=\varnothing,\ R\ne\varnothing &:& \text{append},\\
D\ne\varnothing,\ R=\varnothing &:& \text{prune},\\
D\ne\varnothing,\ R\ne\varnothing &:& \text{exchange}.
\end{array}
\]
If \(\mathcal A\) is the current active support and \(\mathcal P\) is the append
candidate pool, the exact global swap problem is
\[
(D^\star,R^\star)
\in
\argmax_{\substack{D\subseteq\mathcal A\\ R\subseteq\mathcal P}}
\Delta_{\rm patch}(D,R),
\]
with the declared cost and safety constraints enforced.
This subset-search problem is NP-hard in general: it contains best-subset
selection as a limiting case, and the number of possible support patches grows
combinatorially with the active and candidate pools.

\phys{Exchange replaces stale directions with new directions in one checkpoint
edit. The append side still needs append confirmation. The prune side still
needs deletion safety. Exact global swap is therefore a reference optimization
problem, while AP--McLachlan uses scored local batches and guarded commits.}

\section*{26. Ordering of the full algorithmic logic}

\micro{The conceptual order is fixed-support dynamics first, numerical repair
second, structural edits third.}

\deriv
Fixed support means the active generator/coordinate set is held fixed while the
McLachlan velocity is solved. Structural edits are append, prune, and exchange
changes to that active support.
At checkpoint \(k\):
\[
\begin{aligned}
&\ket{\psi_k}=U(\theta_k)\ket{\psi_{\rm ref}},\\
&T_k,\ \ket{b_k},\ G_k,\ f_k,\\
&\dot\theta_k^{(0)}=(G_k+\Lambda_0)_{\rm stab}^{-1}f_k.
\end{aligned}
\]
Then:
\[
\begin{aligned}
&\text{solve repair if }
\rho_{\rm num},\ \Delta t\norm{T\dot\theta},\ \eta_k
\text{ require it},\\
&\text{append if }\rho(T_k,\ket{b_k})\text{ remains structurally large},\\
&\text{prune if grouped loss and history/no-harm checks are safe},\\
&\text{exchange if both sides pass their side constraints}.
\end{aligned}
\]

\phys{This ordering prevents confusing numerical inverse failure with
tangent-span insufficiency. It also prevents prune/exchange notation from obscuring
the simpler append-only and prune-only cases.}

\section*{27. Runtime parameters versus derived objects}

\micro{A threshold is a run setting; a residual, Gram matrix, or velocity is a
checkpoint-derived object.}

\deriv
At a checkpoint,
\[
\ket\psi,\theta,T,\ket b,G,f,\dot\theta
\]
are derived objects. The diagnostic quantities are
\[
\rho(T,\ket b),\qquad \rho_{\rm num}(\ell),\qquad \eta_k.
\]
The runtime inverse settings include
\[
\tau,\qquad \Lambda,\qquad \epsilon.
\]
The runtime decision thresholds include
\[
\tau_{\rm miss},\qquad
\tau_{\rm num},\qquad
\tau_\Delta,\qquad
\tau_{\rm kink}.
\]
The structural choices are
\[
R,\qquad D,\qquad \nu=(D,R).
\]

\phys{The code should record both settings and derived telemetry. Settings say
which rule was declared; derived objects say what happened at the checkpoint.}

\section*{28. Primitive expansion of \(G\dot\theta=f\)}

\micro{The implementation ultimately solves a real matrix equation whose
entries are real parts of complex Hilbert-space overlaps.}

\deriv
Let
\[
T=(T_{aj})_{1\le a\le d,\;1\le j\le n}\in\C^{d\times n},
\qquad
\ket b\leftrightarrow b=(b_a)_{a=1}^{d}\in\C^d.
\]
Then
\[
G=(G_{ij})_{1\le i,j\le n}\in\R^{n\times n},
\qquad
f=(f_i)_{i=1}^{n}\in\R^n,
\]
with
\[
\begin{aligned}
G_{ij}
&=\Real\sum_{a=1}^{d}T_{ai}^*T_{aj},\\
f_i
&=\Real\sum_{a=1}^{d}T_{ai}^*b_a.
\end{aligned}
\]
The velocity equation is
\[
\sum_{j=1}^{n}G_{ij}\dot\theta_j=f_i,
\qquad
1\le i\le n.
\]
Substituting the primitive definitions gives
\[
\sum_{j=1}^{n}
\left(
\Real\sum_{a=1}^{d}T_{ai}^*T_{aj}
\right)
\dot\theta_j
=
\Real\sum_{a=1}^{d}T_{ai}^*b_a.
\]

\phys{This is the base measurement-to-solve path: tangent overlaps fill \(G\),
tangent--drift overlaps fill \(f\), and a real linear solve gives parameter
velocity.}

\section*{29. Primitive expansion of represented velocity}

\micro{The state-space velocity \(T\dot\theta\) is the quantity used by repair
guards, not the coordinate vector alone.}

\deriv
Let
\[
\dot\theta=(\dot\theta_j)_{j=1}^{n}\in\R^n.
\]
The represented velocity is
\[
\ket u=T\dot\theta\in\Hh,
\]
with component form
\[
u_a=\sum_{j=1}^{n}T_{aj}\dot\theta_j,
\qquad
1\le a\le d.
\]
Its squared norm is
\[
\begin{aligned}
\norm{\ket u}^2
&=\sum_{a=1}^{d}u_a^*u_a\\
&=
\sum_{a=1}^{d}
\left(\sum_{i=1}^{n}T_{ai}\dot\theta_i\right)^*
\left(\sum_{j=1}^{n}T_{aj}\dot\theta_j\right)\\
&=
\sum_{i=1}^{n}\sum_{j=1}^{n}
\dot\theta_i
\left(\sum_{a=1}^{d}T_{ai}^*T_{aj}\right)
\dot\theta_j.
\end{aligned}
\]
Since the scalar is real,
\[
\norm{\ket u}^2=\dot\theta^\top G\dot\theta.
\]

\phys{The same \(G\) that appears in the solve also computes the physical size
of the represented velocity. This is why state-motion repair uses
\(\Delta t\norm{T\dot\theta}\).}

\section*{30. Primitive expansion of \(E_\ell\)}

\micro{The realized explained drift is the total drift minus the residual
produced by the specific stabilized velocity.}

\deriv
Let
\[
\dot\theta_\ell=(G+\Lambda_\ell)_{\tau_\ell,\epsilon_\ell}^{-1}f.
\]
The realized residual is
\[
\ket{r_\ell}=T\dot\theta_\ell-\ket b\in\Hh.
\]
In components,
\[
(r_\ell)_a
=
\sum_{j=1}^{n}T_{aj}(\dot\theta_\ell)_j-b_a.
\]
The residual norm is
\[
\begin{aligned}
\norm{r_\ell}^2
&=
\norm{T\dot\theta_\ell-\ket b}^2\\
&=
\dot\theta_\ell^\top G\dot\theta_\ell
-2f^\top\dot\theta_\ell
+\norm{\ket b}^2.
\end{aligned}
\]
Therefore
\[
\begin{aligned}
E_\ell
&=\norm{\ket b}^2-\norm{r_\ell}^2\\
&=
2f^\top\dot\theta_\ell
-\dot\theta_\ell^\top G\dot\theta_\ell.
\end{aligned}
\]

\phys{This is the exact expression Paper II uses to detect whether the
implemented inverse has lost useful tangent motion relative to the supported
least-squares reference.}

\section*{31. Primitive expansion of append block matrices}

\micro{Append introduces a second tangent matrix \(U\), and all augmented
objects are built by ordinary block matrix multiplication.}

\deriv
Let
\[
U=(U_{a\mu})_{1\le a\le d,\;1\le \mu\le r}\in\C^{d\times r}.
\]
Then
\[
T_R=[T~U]\in\C^{d\times(n+r)}.
\]
The augmented Gram matrix is
\[
G_R=\Real(T_R^\dagger T_R)
=
\begin{bmatrix}
G&B\\
B^\top&C
\end{bmatrix},
\]
with
\[
\begin{aligned}
B_{i\mu}
&=\Real\sum_{a=1}^{d}T_{ai}^*U_{a\mu},\\
C_{\mu\nu}
&=\Real\sum_{a=1}^{d}U_{a\mu}^*U_{a\nu}.
\end{aligned}
\]
The augmented force is
\[
f_R=\Real(T_R^\dagger\ket b)
=
\begin{bmatrix}
f\\
q
\end{bmatrix},
\qquad
q_\mu=\Real\sum_{a=1}^{d}U_{a\mu}^*b_a.
\]

\phys{The code path that evaluates a candidate append batch has to compute the
old block \(G,f\), the cross block \(B\), the candidate block \(C\), and the
candidate force \(q\).}

\section*{32. Shape check for append gain}

\micro{The append gain expression is well-defined because every matrix and
vector lives in the augmented coordinate space.}

\deriv
The shapes are
\[
\begin{aligned}
G_R&\in\R^{(n+r)\times(n+r)},\\
f_R&\in\R^{n+r},\\
G_{R,{\rm stab}}^{-1}&\in\R^{(n+r)\times(n+r)}.
\end{aligned}
\]
Thus
\[
G_{R,{\rm stab}}^{-1}f_R\in\R^{n+r},
\]
and
\[
f_R^\top G_{R,{\rm stab}}^{-1}f_R\in\R.
\]
The old score has
\[
G_{\rm stab}^{-1}\in\R^{n\times n},
\qquad
f^\top G_{\rm stab}^{-1}f\in\R.
\]
Therefore
\[
\rho_{\rm gain}(R)
=
\frac{[f_R^\top G_{R,{\rm stab}}^{-1}f_R-f^\top G_{\rm stab}^{-1}f]_+}
{\norm{\ket b}^2}
\]
is a scalar.

\phys{A failed augmented solve should block append admission. The candidate
support is not committed unless the augmented geometry can produce a finite
stabilized velocity.}

\section*{33. Drive profile versus drive-aligned tangent}

\micro{The drive coefficient changes the Hamiltonian; the drive-aligned ansatz
factor changes only the available tangent directions.}

\deriv
Let \(H_0\) be the time-independent Hamiltonian part, \(D=D^\dagger\) the
Hermitian drive operator, and \(c(t)\in\R\) the scalar drive profile. Let the
driven Hamiltonian be
\[
H(t)=H_0+c(t)D,
\qquad
D=D^\dagger.
\]
The physical drift contains
\[
b(t)
=-\ii\bigl(H_0+c(t)D-\langle H(t)\rangle_\psi\Id\bigr)\psi.
\]
A zero-angle ansatz factor is
\[
U_D(\phi)=\exp(-\ii\phi D).
\]
At \(\phi=0\),
\[
U_D(0)=\Id,
\qquad
U_D(0)\psi=\psi.
\]
The added tangent is
\[
\partial_\phi U_D(\phi)\psi\big|_{\phi=0}
=-\ii D\psi,
\qquad
u_D=\Pi_\psi(-\ii D\psi).
\]

\phys{The Hamiltonian drive supplies the physical drift. The zero-angle
drive-aligned ansatz factor supplies a tangent direction that can represent that
drift. It is not counted as the append mechanism being tested in append
diagnostics.}

\section*{34. Per-child versus macro coordinates}

\micro{The same logical generator can define different tangent spaces depending
on how its Pauli-polynomial children are parameterized.}

\deriv
A Pauli-polynomial child is one Pauli term in a decomposed Hermitian generator.
Let
\[
A=\sum_{\mu=1}^{r}c_\mu P_\mu,
\qquad
P_\mu=P_\mu^\dagger.
\]
Per-child coordinates use
\[
U_{\rm child}(\phi)
=
\prod_{\mu=r}^{1}\exp(-\ii\phi_\mu c_\mu P_\mu),
\qquad
\phi\in\R^r.
\]
The zero-angle tangent columns are
\[
u_\mu=\Pi_\psi(-\ii c_\mu P_\mu\psi),
\qquad
1\le\mu\le r.
\]
Macro/shared coordinates use
\[
U_{\rm macro}(\alpha)
=
\prod_{\mu=r}^{1}\exp(-\ii\alpha c_\mu P_\mu),
\qquad
\alpha\in\R.
\]
The zero-angle macro tangent is
\[
u_{\rm macro}
=
\sum_{\mu=1}^{r}u_\mu.
\]

\phys{Per-child mode can append or prune individual Pauli-polynomial children.
Macro mode ties the children together into one shared coordinate. These are
different variational tangent geometries.}

\section*{35. Exact reference boundary}

\micro{Exact trajectories can evaluate the result after the run, but they cannot
enter the AP decision equations.}

\deriv
An exact reference is a separately computed classical trajectory used only for
post-run scoring. The online decision variables are functions of
\[
\psi,\quad H(t),\quad T,\quad b,\quad G,\quad f,\quad U,\quad D.
\]
The exact diagnostic trajectory has the form
\[
\psi_{\rm exact}(t)
=
\mathcal T\exp\left(-\ii\int_0^tH(s)\,\dd s\right)\psi(0).
\]
Reporting errors can use
\[
\Delta O(t)
=
\bra{\psi(t)}O\ket{\psi(t)}
-
\bra{\psi_{\rm exact}(t)}O\ket{\psi_{\rm exact}(t)}.
\]
Let \(C_k\) be the candidate tangent data at the checkpoint and let \(s_k\) be
the runtime setting vector. Controller decisions must have the form
\[
\nu_k
=
\pi\bigl(T_k,b_k,G_k,f_k,C_k,s_k\bigr),
\]
with no \(\psi_{\rm exact}\) argument.

\phys{Energy, doublon, site occupation, and exact overlays are reporting tools.
They do not tune append, prune, repair, or solve choices in a strict AP row.}

\section*{36. Unsupported checkpoint}

\micro{Unsupported means the diagnostic health checks did not all certify the
accepted rung; it is distinct from a runtime crash.}

\deriv
Let \(A\subseteq C\) be the acceptable candidate repair set:
\[
A=\{\ell\in C:\mathcal A(\ell)=1\}.
\]
The predicate \(\mathcal A(\ell)\) equals \(1\) exactly when candidate \(\ell\)
passes the acceptability conditions from Section 12.
If
\[
A\ne\varnothing,
\]
then
\[
\ell^\star\in\argmin_{\ell\in A}J(\ell),
\qquad
\operatorname{unsupported}=0.
\]
If
\[
A=\varnothing
\]
and at least one finite candidate exists, then
\[
\ell^\star\in\argmin_{\ell\in C_{\rm finite}}J(\ell),
\qquad
\operatorname{unsupported}=1.
\]

\phys{The trajectory can continue for scientific diagnostics even when a
checkpoint is unsupported. Strict invalid states, NaN/Inf, impossible
dimensions, and missing required data remain true execution failures.}

\section*{37. Primitive Schur append block}

\micro{The Schur object is built from the same primitive tangent overlaps as
the McLachlan solve.}

\deriv
Let \(U=[u_1~\cdots~u_r]\in\C^{d\times r}\) be the candidate tangent block.
The primitive overlap arrays are
\[
\begin{aligned}
G&=\Real(T^\dagger T)\in\R^{n\times n},\\
B&=\Real(T^\dagger U)\in\R^{n\times r},\\
C&=\Real(U^\dagger U)\in\R^{r\times r},\\
f&=\Real(T^\dagger\ket b)\in\R^n,\\
q&=\Real(U^\dagger\ket b)\in\R^r.
\end{aligned}
\]
Componentwise,
\[
\begin{aligned}
G_{ij}
&=\Real\sum_{a=1}^{d}T_{ai}^\ast T_{aj},
&&1\le i,j\le n,\\
B_{i\mu}
&=\Real\sum_{a=1}^{d}T_{ai}^\ast U_{a\mu},
&&1\le i\le n,\;1\le\mu\le r,\\
C_{\mu\nu}
&=\Real\sum_{a=1}^{d}U_{a\mu}^\ast U_{a\nu},
&&1\le\mu,\nu\le r,\\
f_i
&=\Real\sum_{a=1}^{d}T_{ai}^\ast b_a,
&&1\le i\le n,\\
q_\mu
&=\Real\sum_{a=1}^{d}U_{a\mu}^\ast b_a,
&&1\le\mu\le r.
\end{aligned}
\]

\phys{Every term above is a tangent-overlap or tangent-force quantity. The
append calculation does not need an exact trajectory; it needs the current
state, the current Hamiltonian drift, old tangent columns, and candidate
tangent columns.}

\micro{The Schur subtraction removes the part of the candidate block that the
old tangent span already explains.}

\deriv
Assume \(G^{-1}\) denotes the retained stabilized inverse used for this local
derivation. Then
\[
\begin{aligned}
S&=C-B^\top G^{-1}B\in\R^{r\times r},\\
w&=q-B^\top G^{-1}f\in\R^r.
\end{aligned}
\]
For \(1\le\mu,\nu\le r\),
\[
S_{\mu\nu}
=
C_{\mu\nu}
-
\sum_{i=1}^{n}\sum_{j=1}^{n}
B_{i\mu}(G^{-1})_{ij}B_{j\nu}.
\]
For \(1\le\mu\le r\),
\[
w_\mu
=
q_\mu
-
\sum_{i=1}^{n}\sum_{j=1}^{n}
B_{i\mu}(G^{-1})_{ij}f_j.
\]
The candidate-space solve is
\[
\alpha=S^{-1}w.
\]
The Schur gain is
\[
\Delta_{\rm Schur}
=
w^\top S^{-1}w
=
\sum_{\mu=1}^{r}\sum_{\nu=1}^{r}
w_\mu(S^{-1})_{\mu\nu}w_\nu.
\]

\phys{The candidate block is admitted only for its leftover tangent content.
If \(U\) is already explained by \(T\), then \(S\) becomes small or singular and
the retained novelty check blocks the append.}

\section*{38. Augmented solve confirmation before append}

\micro{The selected append batch is checked by solving the full augmented
McLachlan problem before it is committed.}

\deriv
The augmented tangent matrix is
\[
T_+=
\begin{bmatrix}
T&U
\end{bmatrix}
\in\C^{d\times(n+r)}.
\]
The augmented Gram and force are
\[
G_+
=
\Real(T_+^\dagger T_+)
=
\begin{bmatrix}
G&B\\
B^\top&C
\end{bmatrix},
\qquad
f_+
=
\Real(T_+^\dagger\ket b)
=
\begin{bmatrix}
f\\
q
\end{bmatrix}.
\]
Let \(z=(v,\alpha)\in\R^{n+r}\), where \(v\in\R^n\) are old-support velocity
coordinates and \(\alpha\in\R^r\) are candidate-block velocity coordinates. The
augmented least-squares problem is
\[
z^\star
\in
\argmin_{z\in\R^{n+r}}\norm{T_+z-b}^2.
\]
The stabilized solve used for confirmation is
\[
z^\star_\ell
=
(G_++\Lambda_{+,\ell})_{\tau_\ell,\epsilon_\ell}^{-1}f_+.
\]
The represented augmented velocity is
\[
u_+=T_+z^\star_\ell.
\]

\phys{This is the pre-commit check. The batch is not accepted merely because a
ranking score liked it. The appended support must produce a finite, stabilized,
state-space McLachlan solve at the current checkpoint.}

\micro{The confirmation residual is the realized miss after the candidate is
added at zero amplitude.}

\deriv
The old and augmented realized misses are
\[
\begin{aligned}
R_{\rm old}
&=\norm{T\dot\theta_\ell-b}^2,\\
R_+
&=\norm{T_+z^\star_\ell-b}^2.
\end{aligned}
\]
The confirmed gain fraction is
\[
\gamma_+
=
\frac{[R_{\rm old}-R_+]_+}{\norm{\ket b}^2}.
\]
The append batch may commit only if
\[
\begin{aligned}
z^\star_\ell&\in\R^{n+r},\\
\norm{z^\star_\ell}&<\infty,\\
\norm{u_+}&<\infty,\\
\gamma_+&\ge\tau_{\rm gain},\\
S&\succeq_{\rm retained}\tau_{\rm Schur}.
\end{aligned}
\]

\phys{This closes the gap between candidate scoring and actual propagation.
The exact reference remains outside the decision; the confirmation uses only
the augmented tangent geometry and the current drift.}

\section*{39. Primitive numerical miss chain}

\micro{Numerical miss compares the realized stabilized solve against the best
supported tangent fit.}

\deriv
The stabilized velocity is
\[
\dot\theta_\ell
=
(G+\Lambda_\ell)_{\tau_\ell,\epsilon_\ell}^{-1}f.
\]
The realized explained drift is
\[
E_\ell
=
2f^\top\dot\theta_\ell
-
\dot\theta_\ell^\top G\dot\theta_\ell.
\]
Componentwise,
\[
E_\ell
=
2\sum_{i=1}^{n}f_i\dot\theta_{\ell,i}
-
\sum_{i=1}^{n}\sum_{j=1}^{n}
\dot\theta_{\ell,i}G_{ij}\dot\theta_{\ell,j}.
\]
The realized squared miss is
\[
\epsilon_\ell^2
=
\norm{\ket b}^2-E_\ell.
\]
The supported best-case explained drift is
\[
E_\star
=
f^\top G^+f,
\]
and the supported best-case miss is
\[
\epsilon_\star^2
=
\norm{\ket b}^2-E_\star.
\]
The numerical miss is
\[
\rho_{\rm num}(\ell)
=
\frac{[\epsilon_\ell^2-\epsilon_\star^2]_+}{\norm{\ket b}^2}.
\]

\phys{The same tangent span can be structurally sufficient but numerically
mis-solved. Then the implemented inverse loses drift that the supported
least-squares geometry says was available.}

\micro{The primitive substitution shows exactly where \(T\), \(G\), and \(f\)
enter the realized miss.}

\deriv
\[
\begin{aligned}
\epsilon_\ell^2
&=\norm{T\dot\theta_\ell-b}^2\\
&=\dot\theta_\ell^\top G\dot\theta_\ell
-2f^\top\dot\theta_\ell
+\norm{\ket b}^2\\
&=
\sum_{i,j}\dot\theta_{\ell,i}G_{ij}\dot\theta_{\ell,j}
-2\sum_i f_i\dot\theta_{\ell,i}
+\sum_a b_a^\ast b_a.
\end{aligned}
\]
With
\[
\begin{aligned}
G_{ij}&=\Real\sum_a T_{ai}^\ast T_{aj},\\
f_i&=\Real\sum_a T_{ai}^\ast b_a,
\end{aligned}
\]
the fully substituted scalar is
\[
\begin{aligned}
\epsilon_\ell^2
={}&
\sum_{i,j}\dot\theta_{\ell,i}\dot\theta_{\ell,j}
\Real\sum_a T_{ai}^\ast T_{aj}\\
&-
2\sum_i\dot\theta_{\ell,i}
\Real\sum_a T_{ai}^\ast b_a
+
\sum_a b_a^\ast b_a.
\end{aligned}
\]

\phys{This is why diagnostics should use \(T\dot\theta\) in state space. The
coordinate vector only matters through the state-space velocity it synthesizes.}

\section*{40. Why damping can help or hurt}

\micro{Regularization changes the inverse mode-by-mode, so it can cure an
unstable solve or over-damp a useful one.}

\deriv
Let
\[
G=Z\diag(s_1,\ldots,s_n)Z^\top,
\qquad
Z^\top Z=\Id.
\]
The columns \(z_\alpha\) of \(Z\) are orthonormal eigenmodes of \(G\), and
\(s_\alpha\) are their metric strengths.
Let
\[
a=Z^\top f.
\]
Thus \(a_\alpha=z_\alpha^\top f\) is the force component along eigenmode
\(z_\alpha\).
For a scalar ridge \(\lambda_\ell\) and solve damping \(\epsilon_\ell\), the
retained inverse gives
\[
\dot\theta_\ell
=
\sum_{\alpha=1}^{n}
\chi_{\ell,\alpha}
\frac{a_\alpha}{s_\alpha+\lambda_\ell+\epsilon_\ell}
z_\alpha.
\]
The retained-mode contribution to represented velocity is
\[
u_\ell
=
T\dot\theta_\ell
=
\sum_{\alpha=1}^{n}
\chi_{\ell,\alpha}
\frac{a_\alpha}{s_\alpha+\lambda_\ell+\epsilon_\ell}
Tz_\alpha.
\]
For a retained mode,
\[
\frac{\partial}{\partial\lambda_\ell}
\left(
\frac{a_\alpha}{s_\alpha+\lambda_\ell+\epsilon_\ell}
\right)
=
-
\frac{a_\alpha}{(s_\alpha+\lambda_\ell+\epsilon_\ell)^2}.
\]

\phys{Increasing ridge damps every retained coordinate mixture. This can remove
a state-space jump caused by small unstable modes, but it can also suppress
real drift if the mode was useful. That is why \(\rho_{\rm num}\) triggers a
candidate search rather than dictating one direction.}

\micro{A high condition number is a response-risk signal, not the main
noiseless trigger.}

\deriv
For retained eigenmodes
\[
I_\ell=\{\alpha:\chi_{\ell,\alpha}=1\},
\]
define
\[
\kappa_{\rm eff}(\ell)
=
\frac{\max_{\alpha\in I_\ell}s_\alpha}
{\min_{\alpha\in I_\ell}s_\alpha}.
\]
The least-intervention score may contain
\[
w_\kappa\phi_\kappa(\ell),
\]
but repair entry in the noiseless diagnostic route comes from
\[
g_\rho\vee g_\Delta\vee g_t\vee g_{\rm inv}.
\]

\phys{A large \(\kappa_{\rm eff}\) makes a candidate less attractive because
small errors can be amplified. It does not by itself force the controller into
global repair at every checkpoint.}

\section*{41. Local subdivision as the first state-motion cure}

\micro{When the vector field is plausible but the step is too large,
subdivision reduces integration error without changing the tangent manifold.}

\deriv
Subdivision means replacing one interval of length \(\Delta t\) by \(r\)
shorter intervals while using the same local vector-field construction.
For a first-order local step,
\[
\psi(t+\Delta t)
=
\psi(t)+\Delta t\,u(t)+O(\Delta t^2).
\]
With \(r\) substeps of size \(\delta t=\Delta t/r\),
\[
\psi(t+\Delta t)
=
\psi(t)
+
\sum_{m=0}^{r-1}\delta t\,u(t_m)
+
O(r\delta t^2).
\]
Since
\[
r\delta t^2
=
r\left(\frac{\Delta t}{r}\right)^2
=
\frac{\Delta t^2}{r},
\]
the leading local discretization scale decreases as \(r^{-1}\).

\phys{If the defect is a temporal kink or an excessive state-space step, the
first cure is to resolve the time interval more finely. Heavy damping changes
the vector field; subdivision first asks whether the same vector field was
being stepped too coarsely.}

\section*{42. Support change continuity}

\micro{Appending at zero amplitude changes the tangent space before it changes
the state.}

\deriv
State-continuous means the state vector is unchanged at insertion.
Tangent-discontinuous means the set of available first-order velocities changes
at the same state. Let \(A=A^\dagger:\Hh\to\Hh\) and let
\[
\psi_+(\theta,\phi)
=
\exp(-\ii\phi A)\psi(\theta).
\]
At \(\phi=0\),
\[
\psi_+(\theta,0)=\psi(\theta).
\]
The new tangent is
\[
\left.\frac{\partial}{\partial\phi}\psi_+(\theta,\phi)\right|_{\phi=0}
=
-\ii A\psi(\theta).
\]
After horizontal projection,
\[
u_A
=
\Pi_\psi(-\ii A\psi).
\]
Thus
\[
T_+
=
\begin{bmatrix}
T&u_A
\end{bmatrix}.
\]

\phys{A zero-angle append is state-continuous but tangent-discontinuous. The
state does not jump, while the available McLachlan velocities can change
immediately at the append checkpoint.}

\micro{Pure append, pure prune, and exchange are the same patch object with
different empty sides.}

\deriv
Let \(D\) be a deletion set and \(R\) an append set. The patched tangent matrix
is
\[
T_\nu
=
\begin{bmatrix}
T_{-D}&U_R
\end{bmatrix}.
\]
The limiting cases are
\[
\begin{array}{ccl}
D=\varnothing,\ R=\varnothing &\Longrightarrow& T_\nu=T,\\[1mm]
D=\varnothing,\ R\ne\varnothing &\Longrightarrow& T_\nu=[T~U_R],\\[1mm]
D\ne\varnothing,\ R=\varnothing &\Longrightarrow& T_\nu=T_{-D},\\[1mm]
D\ne\varnothing,\ R\ne\varnothing &\Longrightarrow& T_\nu=[T_{-D}~U_R].
\end{array}
\]

\phys{Append and prune are easiest to learn separately, but Paper II's full
object is the support patch. The same least-squares diagnostics decide whether
the patch reduces structural tangent miss without creating unacceptable numerical
or continuity risk.}

\section*{43. A two-coordinate null-direction example}

\micro{A toy tangent matrix makes non-identifiability visible without adding
physics assumptions.}

\deriv
Let
\[
T=
\begin{bmatrix}
1&1\\
0&0
\end{bmatrix},
\qquad
a=
\begin{bmatrix}
1\\
-1
\end{bmatrix}.
\]
Then
\[
Ta=
\begin{bmatrix}
1&1\\
0&0
\end{bmatrix}
\begin{bmatrix}
1\\
-1
\end{bmatrix}
=
\begin{bmatrix}
0\\
0
\end{bmatrix}.
\]
The Gram matrix is
\[
G=T^\top T
=
\begin{bmatrix}
1&1\\
1&1
\end{bmatrix},
\]
and therefore
\[
Ga=
\begin{bmatrix}
1&1\\
1&1
\end{bmatrix}
\begin{bmatrix}
1\\
-1
\end{bmatrix}
=
\begin{bmatrix}
0\\
0
\end{bmatrix}.
\]
Only the sum direction is identifiable:
\[
T\dot\theta
=
\begin{bmatrix}
\dot\theta_1+\dot\theta_2\\
0
\end{bmatrix}.
\]
The difference direction \(\dot\theta_1-\dot\theta_2\) is invisible to the
tangent map.

\phys{The two coordinates look distinct in parameter space, but locally they
draw the same state-space arrow. McLachlan should not spend inverse weight on
the invisible difference direction.}

\section*{44. Compact mental model}

\micro{The AP--McLachlan route is one repeated least-squares geometry with
three different questions.}

\deriv
\[
(T,\ket b)
\longmapsto
(G,f,\norm{\ket b}^2)
\longmapsto
G^+
\longmapsto
\dot\theta
\longmapsto
T\dot\theta.
\]
\[
\dot\theta
=
\argmin_v\norm{Tv-b}^2.
\]
\[
\begin{aligned}
\dot\theta_\ell
&=(G+\Lambda_\ell)_{\tau_\ell,\epsilon_\ell}^{-1}f,\\
u_\ell&=T\dot\theta_\ell.
\end{aligned}
\]
\[
\begin{aligned}
T_{\rm app}&=[T~U_R],\\
\Delta_{\rm app}
&=E(T_{\rm app})-E(T).
\end{aligned}
\]
\[
\begin{aligned}
T_{\rm pr}&=T_{-D},\\
\Delta_{\rm pr}
&=E(T)-E(T_{\rm pr}).
\end{aligned}
\]
\[
T_\nu=[T_{-D}~U_R].
\]

\phys{All major formulas come from the same object: least-squares projection of
the Schrödinger drift \(b\) onto a tangent span. The inverse determines the
velocity; repair stabilizes that inverse; append expands the span; prune
contracts the span only after safety checks.}

\end{document}
